\scanpage{359}
% @ head of sheet
\mspage{83}
% @ top, lines 1–5
Or $i\in I_\infty$ signifie que tous les $\Omega^n(X_i,x_i)$ sont $0$-connexes, ce qui signifie aussi que $X_i$ est $\mathbb W_\infty$-asphérique (puisqu'il est déjà $0$-connexe) - (cf $\overline{\underline{\mathrm{XII}}}$ \unsure{cor.} p. 190). Cela achève la démonstration du th. 5.4.

% @ upper third, section heading
\deleted{\unsure{5}}6. \uline{\deleted{\unsure{L'axiome de descente} $\mathrm{Loc}(\unsure{10})$}} \inserted{\uline{La condition} $\mathrm{Loc}(5,D_{W_{\unsure{\omega|W|w}}})$ \uline{implique} $W\subset\mathbb W_\infty$.}

% @ centre, introduction and theorem
Je vais redonner ici, pour mémoire, la démonstration du th. 1 de $\overline{\underline{\mathrm{XII}}}$ :

\uline{Théorème} 6.1. Soit $W$ un localiseur fondamental satisfaisant à $D_{W_\omega}$ (où $W_\omega$ est le plus petit de tous les localiseurs fondamentaux). Alors on a $W=\mathrm{Fl}(\mathrm{Cat})$, \uline{ou} $W\subset\mathbb W_\infty$ (suivant que $W$ ne satisfait pas à l'axiome de connexité $\mathrm{Loc}(5)$, ou qu'il y satisfait).

% @ lower third, references following the theorem
Je vais faire la démonstration en supposant que $W$ satisfait à $\mathrm{Loc}(5)$. Pour la démonstration du fait que dans le cas contraire, on a $W=\mathrm{Fl}(\mathrm{Cat})$, je renvoie à $\overline{\underline{\mathrm{XII}}}$ th. 2 p. 20. (Voir aussi le résultat \inserted{plus général et plus faible,} très voisin, \deleted{\unsure{p. 12}} \inserted{rappelé \unsure{en}} prop 2.7. page 12.)

% @ foot, Lemma 6.2
\uline{Lemme} 6.2. Soient $W,W'$ deux loc fond., \inserted{$W$} satisfaisant à $\mathrm{Loc}(7\,D_{W_\omega})$. Pour qu'on ait $W\subset W'$, il suffit qu'on ait $\mathrm{Asph}_W\subset\mathrm{Asph}_{\unsure{W_\infty|W'}}$
% archive sheet number: 358
